\chapter{State of the Art} \label{chap:sota}

In this chapter, a brief review of the state of the art regarding whispered speech conversion  is provided, focusing on the techniques that have been previously employed to tackle the problem and their respective results and conclusions. Two main categories of whispered to normal voice conversion techniques exist: rule-based and machine learning-based approaches \cite{perrotinGlottalFlowSynthesis2020}.

\section{Rule-based}

Firstly, rule-based approaches work by modeling specific aspects of speech production in order to add voicing to the whispered signals, by defining deterministic transformation rules. These can be obtained either by empirical analysis or statistical data regarding the acoustical differences between voiced and unvoiced speech, although the simplicity of these models has an impact in both speech naturalness and quality \cite{lianWhisperNormalSpeech2019}. This can be done in several ways, as shown further in this section, but this is mostly implemented by identifying and segmenting the input whispered speech according to whether they represent potentially voiced or unvoiced segments and then adding said voicing components.

As pointed out by Perrotin et.al. \cite{perrotinGlottalFlowSynthesis2020}, these frameworks usually first decompose whisper, using linear prediction in order to estimate a parametric filter representing the spectral speech envelope, and then inverting said filter to obtain and excitation signal. Modification of the filter parameters is done afterwards, while also determining the vocalized source excitation parameters, which are then used to excite said filter and thus generate the enhanced speech.

\subsection{MELP}

One relevant work that is important to point out is that of Morris and Clement \cite{morrisReconstructionSpeechWhispers2002a}, in which a modified MELP model is employed. This is a standard way of modeling speech production, utilizing LPC to model the vocal tract, as was mentioned in Chapter \ref{chap:back}. MELP is based on mixed excitation, meaning that both periodic pulses and noise are used as excitation signals for the voiced and unvoiced/whispered sections, respectively. Both components are then mixed dynamically according to what is often referred to in literature as the aperiodicity coefficient \cite{perrotinGlottalFlowSynthesis2020}, resulting from the analysis of each signal segment, in order to produce a more natural sound. 

In \cite{morrisReconstructionSpeechWhispers2002a}, the MELP model is built upon by injecting a synthesized periodic wave to the source excitation signal, which in turn adds voicing to whispered signal, although no mention is made about the pulse shape that is used. In order to provide user feedback control, the fundamental frequency of this synthesized signal is then controlled by the sound intensity. The filter part of MELP is first modified through the implementation of a long-term, minimum phase filter, which is intended to compensate for the differences between the voiced and unvoiced signal spectrums. After the initial filtering, a correction of the formant frequencies and their bandwidth is performed, shifting them and narrowing their bands. This method does not perform any speech segmentation and instead applies permanent voicing to the low-end frequencies up to 3kHz and aperiodic excitation to the higher frequencies up to 4kHz. 

The results do show that the quality of the voice signal predicted from the whispered one is good and that the method works reasonably well, and further research \cite{huangReconstructionWhisperChinese2012} has also shown that good intelligibility can be achieved with this method. Yet, as pointed out in \cite{mcloughlinReconstructionContinuousVoiced2013,ferreiraImplantationVoicingWhispered2016}, the system is not suitable to operate in real-time applications, depending on the smoothing filter that is utilized. Furthermore, this solution does require original voice samples to be available for generating the excitation for the MELP analysis, which are not available in most practical cases.

\subsection{CELP}

Another speech codec based alternative is proposed by McLoughlin et.al. \cite{mcloughlinReconstructionPhonatedSpeech2015} using excitation codebooks with Code Excited Linear Prediction (CELP). This codec operates by first doing short-time predictions with an LPC model, thus modelling the spectral content of the speech segment, followed by an Long-term Prediction (LTP) which in turn detects the pitch. The excitations are then determined from a codebook of random Gaussian sequences using an analysis-by-synthesis approach to minimize the residual error.

This paper describes how the model is expanded upon, first by implementing formant modification and artificial glottal modulation. The formant frequencies and their amplitudes are estimated in a novel way to avoid implausible formants inferred from the noise-like nature of whisper. Preprocessing is applied to the whispered signal, firstly by identifying wether the current segment under analysis is unvoiced or not through the use of a modified G.729 detector \cite{beritelliPerformanceEvaluationComparison2002}, using both signal power and zero crossing rate to determine this. This stage is followed by a whispered phoneme classifier that analyses the spectrum for each speech segment comparing the power distribution across different bands, resulting in a classification into either plosives, fricatives or vowels. This preprocessing is then followed by a spectral enhancement stage that is applied to the whispered segments that are classified into vowels, as these are the parts where any formant imperfections do become apparent to listeners \cite{klattAnalysisSynthesisPerception1990}.



\subsection{Glottal Flow synthesis}

Perrotin et.al \cite{perrotinGlottalFlowSynthesis2020} propose a new parametric approach for converting whispered speech into natural sounding voiced speech, through the use of a sophisticated deterministic model that attempts to accurately parameterize both the glottal pulse excitation and the vocal tract. This is slightly different in function from the previously presented solutions as the authors claim that this approach is voice-based, meaning that the source model does not consider only the excitation but also the glottis filter and the lip radiation. This is opposed to the previous methods which would fall into the DSP-based category of modelling, where the source is just the excitation while the filtering is a result of the glottis, VT and the lips.

The proposed method synthesizes an excitation source based on a pulse train combination of harmonic frequencies, whose frequencies are based on a fundamental frequency multiplied by a random factor above a certain $F_{VUV}$ frequency. The phase is also slightly randomized for each harmonic as to introduce natural sounding roughness. This is done after the whispered signal goes through a preprocessing stage, where inverse filtering is done to first obtain both vocal tract and glottis filter coefficients, in this case using GFM-IAIF \cite{alkuGlottalWaveAnalysis1992} to estimate both from the whispered signal. A new glottis filter is then synthesized based on configurable parameters which allow the user to indicate the range of effort and tension \cite{feugereCantorDigitalisChironomic2017} for the synthesized signal. The synthesized excitation is then filtered through the newly synthesized filter and mixed with the original glottal flow derivative according to whether the current segment is voiced or unvoiced. Finally, the original VT filter filters these segments. Another important aspect that is introduced in this work is the use of the spectral centre of gravity in order to determine the voiced/unvoiced decision, which is then used to determine the amount of voicing and noise that is added to the signal, an idea that also inspires our work, as detailed in Chapter \ref{chap:overlapAdd}.

The authors then go on to evaluate this method through subjective and objective tests to measure the naturalness of synthesized speech and the accuracy of voiced/unvoiced (V/UV) decisions, using a MUSHRA \cite{internationaltelecommunicationunionBS15343Method2015} perceptual testing environment. The subjective experiment involved comparing the proposed glottal source model against the baseline pulse train method and natural speech sources using the wTIMIT database, which contains pairs of phonetically balanced sentences. The new glottal source significantly outperformed the baseline pulse train method in terms of naturalness, receiving higher median scores. The accuracy of V/UV decisions revealed a total V/UV error rate of 17.7\%, with a significant negative correlation between the UV-to-V error rate and naturalness scores. Sentences with fewer unvoiced consonants received higher scores, indicating sensitivity to V/UV decisions.


\subsection{Other Speech Conversion Approaches}

Toda et. al. \cite{todaStatisticalVoiceConversion2012} explore advanced methods to enhance body-conducted unvoiced speech for silent speech communication using a nonaudible murmur (NAM) microphone. This type of microphone captures very soft unvoiced speech, such as NAM or whispered voices, while keeping emitted sounds almost inaudible to those around the speaker. The main challenge addressed in the article is the unnatural and less intelligible quality of body-conducted unvoiced speech due to acoustic changes caused by body conduction.

To tackle this issue, the authors propose two primary voice conversion methods. The first method, NAM-to-Speech, aims to convert NAM to normal speech to improve intelligibility. However, this method often results in degraded naturalness because it is challenging to estimate natural fundamental frequency contours from unvoiced speech. The second method, NAM-to-Whisper, converts NAM to a whispered voice, which is a familiar form of unvoiced speech and thus easier to convert accurately. This method significantly enhances both the intelligibility and naturalness of the converted speech compared to NAM-to-Speech. The conversion process relies on Gaussian mixture models (GMMs) to probabilistically transform the acoustic features of body-conducted unvoiced speech into those of natural voices.

Experimental evaluations demonstrate the effectiveness of these proposed methods. NAM-to-Whisper outperforms NAM-to-Speech in both intelligibility and naturalness, making it a more promising solution for enhancing body-conducted unvoiced speech. Additionally, the methods are extended to handle both NAM and body-conducted whispered voice (BCW), improving the versatility and robustness of the system. The article underscores the potential of these statistical voice conversion techniques to overcome the inherent problems of silent speech communication, paving the way for more universal and private speech communication methods.

%Non-audible nurmurs\dots
%
%Statistical Voice Conversion Techniques for Body-Conducted Unvoiced Speech Enhancement \cite{todaStatisticalVoiceConversion2012}
%
%Some have also studied using ingressive speech in order to \dots

\section{Machine Learning Approaches}

As opposed to rule-based, ML approaches rely on learning algorithms to reconstruct voiced speech from whispers. These algorithms are trained on matched data points of natural and whispered speech, requiring a specially curated dataset although some approaches avoid this necessity through non-parallel training models. Recent advancements have introduced a variety of machine learning frameworks to tackle the challenges imposed by WSC, namely the intricate task of estimating the fundamental frequency $F_0$ contour \cite{oliveiraMachineLearningApproaches2022}, and the automated segmentation of speech, both crucial for the prosody characteristics of natural speech. 

Examples of ML approaches in this field are mostly based on Deep Neural Network (DNN) architectures, including Generative Adversarial Networks, Autoencoder, Bidirectional Long Short-Term Memory and Transformer frameworks. One main issue of most of the papers analyzed in \cite{oliveiraMachineLearningApproaches2022} is that most apporaches did not do any sort of perceptual testing and opt to utilize other objective metrics to evaluate the performance of their systems, such as utilizing speech recognition technologies for this purpose. In the following sections, a non-exhaustive review of the most relevant works in this field is presented. It is important to note that the field of WSC can be seen as a niche within a niche in research, as the amount of deep learning audio research, although growing, is still relatively small when compared to other fields such as image processing or natural language processing.

\subsection{Long Short-Term Memory}

In \cite{lianWhisperNormalSpeech2019}, a sequence-to-sequence model is proposed for whispered speech conversion. The Seq2Seq model used is based on a LSTM architecture \cite{cho2014learningphraserepresentationsusing}, which was a pioneering architecture that has been utilized successfully in many natural language processing applications. This proposed model is trained on a dataset of whispered and normal speech pairs, and is capable of generating natural sounding speech from whispered input. The model is composed of two main components, an encoder and a decoder, which are both based on Long Short-Term Memory (LSTM) networks. The encoder is responsible for encoding the input whispered speech, while the decoder generates the normal speech from the encoded representation. The model is trained using a combination of the mean squared error loss and the mel-cepstral distortion loss, which is a measure of the difference between the mel-cepstral coefficients of the generated and target speech.

Experimental results validate the effectiveness of the proposed method which exhibits superior performance across different metrics, such as Mel Cepstral Distance (CD) and Short-Time objective Intelligibility (STOI) when compared to other state-of-the-art methods. The authors also conducted a subjective evaluation of the generated speech, which showed that the proposed method outperformed other methods such as Gaussian Mixture Models (GMM) and Bidirectional LSTMs \cite{todaNAMtospeechConversionGaussian2005, meenakshiWhisperedSpeechNeutral2018} in terms of naturalness and intelligibility.

\subsection{Generative Adversarial Networks}

Parmar et.al. \cite{parmarEffectivenessCrossdomainArchitectures2019} investigated the application of Generative Adversarial Networks (GANs) to model the intricate differences between whispered and normal speech. Two main architectures are focused on, the first being a CycleGAN \cite{zhu2020unpairedimagetoimagetranslationusing} based model and the second being a modified DiscoGAN \cite{shahNovelMMSEDiscoGAN2018}, which are both trained on a dataset of whispered and normal speech pairs. Both models aim to learn the cross-domain relationships, i.e. the mapping between whispered and normal speech enhancing the quality and naturalness.

The authors conducted comprehensive experiments using the whispered TIMIT (wTIMIT) database, training and evaluating their models on parallel datasets of whispered and normal speech. Objective measures such as Mel Cepstral Distortion (MCD) and Root Mean Square Error (RMSE) of log(F0) were used to assess the performance of the conversion systems. The results demonstrated that both DiscoGAN and CycleGAN significantly outperformed traditional GAN models in reducing MCD and RMSE. Interestingly, while DiscoGAN yielded better objective results, CycleGAN was preferred in subjective evaluations for its naturalness, where this preference is attributed to the differences in their loss functions. Findings indicate that these models can produce more natural-sounding normal speech from whispered input, with CycleGAN being particularly favored for its subjective performance. The authors suggest that further improvements can be achieved by exploring high-quality vocoders such as WaveNet \cite{oord2017parallelwavenetfasthighfidelity} to enhance the voice quality of the converted speech.

Another GAN based architecture as been deployed by Patel et. al. \cite{patelCinCGANEffectiveF02021} to improve the conversion of whispered speech to normal speech by predicting the fundamental frequency (F0). The authors introduce the Cycle-in-Cycle GAN (CinC-GAN), an advanced version of CycleGAN, designed to enhance the accuracy of F0 prediction without compromising the mapping of Mel Cepstral Coefficients (MCC). This is crucial because traditional CycleGAN methods, which rely on separate models for MCC mapping and F0 prediction, often introduce significant non-linear noise in the predicted F0 due to the non-linear DNN layers at the network's output.

This CinC-GAN architecture introduces a joint training methodology that concurrently trains for MCC mapping and F0 prediction. This holistic approach has been shown to reduce the non-linear noise that typically affects the output of sequentially trained models. Finally the paper reports that CinC-GAN significantly outperforms the baseline CycleGAN in both objective and subjective evaluations. Objective metrics, such as Mel Cepstral Distortion (MCD) and Root Mean Square Error (RMSE) of log(F0) also show substantial improvements. Additionally, subjective evaluations reveal that CinC-GAN delivers superior naturalness scores across both seen and unseen speakers, underscoring its effectiveness in generating more natural-sounding speech.z 

\subsection{Transformer Networks}

Transformer networks have more recently been widely used for many natural language processing applications, as their attention mechanisms allow for more efficient handling of sequential data, such as text information or audio signals, on sequence-to-sequence tasks when compared to traditional RNNs (Recurrent Neural Networks) and CNNs (Convolutional Neural Networks). First introduced in the paper by \cite{vaswani2023attention}, the key components of the transformer are it's attention mechanism, which allow the networks to weigh different input tokens in a sequence when encoding each token, and the positional encoding, which allows the network to understand the order of the tokens in the sequence.

Nirajan et al. \cite{niranjanEndtoEndWhisperNatural2021} propose an {\it end-to-end} approach to WSC by adapting a conventional transformer network. This method incorporates three key components, first the standard encoder and decoder found in transformer architectures, along with an additional encoder designed to predict triphones, which are sequences of three phonemes. Their implementation involves two distinct networks; one processes Mel-frequency Cepstral Coefficients (MFCCs), and the other processes smooth spectral features. Although no perceptual tests were conducted, the findings indicate a generally lower word error rate in automatic speech recognition (ASR) systems when using the MFCC-based network.

On the other hand, more recent work conducted by Rekimoto \cite{rekimotoWESPERZeroshotRealtime2023} takes advantage of transformer networks in order to perform speech conversion in what is claimed to be real-time. The system architecture proposed consists on a Speech-to-Unit encoder network that trained on both whispered and normal speech, and it is essentially a modified version of BERT that learns to map both types of data samples to common speech units in a self-supervised manner. Then a separate decoder network based on a simplified version of Microsoft's Fastspeech II \cite{ren2022fastspeech2fasthighquality} handles the conversion of the individual units to mel-spectrograms, which are then converted back to audio using Hifi-GAN \cite{kong2020hifigangenerativeadversarialnetworks}. 

The results of the WESPER system, as detailed in the study, indicate a significant improvement in the perceptual quality of whispered speech when converted to normal speech. Evaluations using Mean Opinion Score (MOS) and MUSHRA \cite{internationaltelecommunicationunionBS15343Method2015} methodologies demonstrated that the converted voices scored higher than whispered voices in terms of quality, naturalness, and prosody preservation. Specifically, the converted speech showed statistically significant improvements in perceived naturalness and reduced hoarseness. Additionally, for individuals with Vocal Fold Polyps (VFP), Spasmodic Dysphonia (SD), and hearing impairments, the converted speech exhibited enhanced intelligibility and naturalness. The study concludes that WESPER offers an arguably real-time solution for effectively converting whispered speech, without compromising on the natural prosody of the original speech.



%\section{Research at FEUP}


